\section{Background and literature review}
\label{sec:bg}

This chapter supplies the frame in which the choices of \cref{sec:robot,sec:rl} are positionned. 
It also fixes the vocabulary, because ``stable'' carries different meanings in this field. 
It present the historical context of bipedal locomotion and sorts theses approaches on seven analytical axes. 
Finally, it narrows to the reinforcement-learning lineage this architecture belongs to.

\subsection{Definitions}
\label{sec:bg-defs}

Each historical family attempts bipedal locomotion by a different strategy. This section fixes the
vocabulary those families share, then sets out the four stability concepts the rest of the chapter
sorts them by.


A few terms recur :
\textbf{Underactuation} is the excess of degrees of freedom over actuators. 
\textbf{Unilateral contact} means a foot can push but not pull, and only inside the \textbf{friction cone}. 
The point where the push applies is the \textbf{centre of pressure}, confined to the contact patch. 
\textbf{Duty factor} is the fraction of a stride each foot spends in stance, and a \textbf{flight phase} is the interval with no foot down. 
The \textbf{sim-to-real gap} is the performance loss on transfer from a simulated plant to the physical one.
In the context of this thesis, \textbf{Failure} is episode termination on body contact, or on leaving the measured workspace. 
\textbf{Stability margin} is reported as a survivor count and a seed variance rather than as a scalar reward, for the reason given in \cref{sec:margin-analysis}.
The reference timescale is the free-topple time of the inverted pendulum, $\sqrt{h/g} \approx 310$\,ms for this robot.
\textbf{Stability} is the ability to continue locomotion when subjected to disturbances. 

\textbf{Static stability} holds when the ground projection of the centre of mass stays inside the
support polygon. It assumes quasi-static motion, flat feet and flat ground. The flat-ground
assumption is the fragile one, because on uneven terrain the support-polygon test can be arbitrarily
wrong~\cite{bretl2008static}. Static stability therefore admits only gaits slow enough for inertia to
be negligible, with at least one foot down throughout.

\textbf{Zero moment point stability} holds when the zero moment point stays inside the support
polygon~\cite{vukobratovic1969}. The criterion releases the quasi-static assumption, so the centre of
mass may travel beyond the support polygon, and it carried the field through
ASIMO~\cite{hirai1998honda}. It still assumes a contact patch of finite area, and still forbids a
flight phase. The test also stays conservative, because it reduces the dynamics to a force balance
and so admits slow motions rather than fully dynamic ones~\cite{chevallereau2003rabbit}. For a point
foot the support polygon degenerates to a point, and the test then never passes while the robot
moves. DASH-01 carries a spherical toe, so neither this nor the previous concepts applies to it.

\textbf{Orbital stability} holds when the spectral radius of the linearised Poincar\'e map falls below one.
If the point is outside the stable radius, the gait might work but is unstable to any tiny perturbation. 
It certifies a whole period of the gait instead of a single instant. 
A flight phase is admissible. It work only by assuming an
exactly known model and a fixed periodic gait, and each new gait need one offline
optimisation~\cite{westervelt2003hzd}. The map is linearised, so the certificate covers small
perturbations only. RABBIT is the point-foot biped the framework was built for~\cite{chevallereau2003rabbit}.

\textbf{Capturability} holds when a reachable foothold exists that brings the reduced model to a stop~\cite{pratt2006capture,koolen2012capturability}. 
Unlike the three previous concepts, which care about whether the present state is valid, what capturability care about is whether a
stopping option exists in the future. It need a reduced model and reachable footholds.

The field has moved down the order of the previous list over the last fifty years, releasing constraints to buy performance~\cite{wieber2016handbook,pratt2006margins,hobbelen2007lcw}.

\subsection{What makes bipedal locomotion hard}
\label{sec:bg-hard}

A legged robot touches the world only through intermittent unilateral contacts. Three consequences
follow.

\textbf{It is underactuated.} Six leg motors against twelve degrees of freedom leaves six base
degrees with no motor attached. The spring angle are not considered as DoF. Those six can be influenced only through the ground. A flight phase
has no ground, so the centre of mass follows a ballistic arc no actuation can alter, and the robot
chooses only the posture it will land in.

\textbf{The inputs are unilateral and saturating.} The centre of pressure moves only within the
contact patch, and since the robot has a point foot, the patch is a point. That leaves foot placement as the
only input. It is chosen once per step, and bounded by the leg's reachable workspace.

\textbf{The timescales are short.} An inverted pendulum of leg length $h$ diverges with time
constant $\sqrt{h/g}$, about 310\,ms in this case, and running contacts last tens of milliseconds. A
correction are time critical, which makes actuator bandwidth a major constraint since we can only react trough foot placement. \Cref{sec:plant-id-drive} measures
this robot's drive at a 0.8\,Hz pole in servo mode and 12 \,Hz in MIT mode. Attentive reader will notice the traquenard that will come back later.

\subsection{Historical context}
\label{sec:bg-history}

As explained in the Definitions section, the families below are sorted by where they get their stability from.

\textbf{Geometry.} WABOT-1 (1973) moved slowly enough for its own momentum to be
neglected~\cite{lim2007waseda}. The dynamic generalisation is the zero moment
point~\cite{vukobratovic1969,vukobratovic2004}, which carried the field for three decades through
Honda's P2 and ASIMO~\cite{hirai1998honda}. The linear inverted pendulum model made it cheap to compute~\cite{kajita2001lipm,kajita2003preview}.

\textbf{The mechanism.} McGeer's passive dynamic walkers descend a ramp passively using gravity alone ~\cite{mcgeer1990}. 
The gait is a limit cycle set by linkage geometry and mass distribution. 
Collins et al.\ added a motor to walk on a flat ground, 
and measured a cost of transport near 0.20, which is similar to humans~\cite{collins2005}.

\textbf{Foot placement.} Raibert's hoppers reframed the problem: a running robot has to be balanced
only during the stance phase rather than continuously~\cite{raibert1986}. 
Control decomposes into thrust at bottom of stance, foot placement about a neutral point, and hip torque applied during stance only.

\begin{equation}
  x_f \;=\; \tfrac{1}{2}\,\dot x\,T_s \;+\; k\,(\dot x - \dot x_d),
  \label{eq:neutral-point}
\end{equation}

$x_f$ is the fore-aft foot placement, measured from the hip at touchdown, in metres. \\
$\dot x$ is the current forward speed in metres per second. \\
$\dot x_d$ is the commanded forward speed, in metres per second.  \\
$T_s$ is the stance duration in seconds. \\
$k$ is a velocity feedback gain, also in seconds.  \\
This equation give the foot placement that will bring the robot back to the commanded speed in one step. 
The first term is the neutral point, which yield no net fore-aft impulse. 
The second term corrects for the speed : running faster than commanded places the foot further forward, so stance decelerates the body.
The lineage runs through BigDog~\cite{raibert2008bigdog} to Atlas.
This equation made the early Boston dynamics possible.

\textbf{Re-solving online. (MPC)} Wieber recycled zero-moment-point walking as model-predictive
control~\cite{wieber2006mpc} by doing on line. Kuindersma et al.\ built the whole-body version that ran Atlas at the
DARPA Robotics Challenge~\cite{kuindersma2016atlas}.

\textbf{Reinforcement learning.} A neural network learn a policy by reinforcement learning to
maximise a reward over millions of simulated episodes.
PPO~\cite{schulman2017ppo} became the workhorse, MuJoCo~\cite{todorov2012mujoco} made it computationally affordable, and
massively parallel simulation cut wall-clock further as demonstrated by Rudin et al. by training a walking policy under 4 minutes~\cite{rudin2021}. 
What used to be stability margin now show up in the form of a sim2real gap.
Previous approach yielded a fixed mathematical law that matched a specific plant (aka robot)
The RL approach allow to train against a multitude of variation of the plant to make it robust to unknown at the expense of some peak performance and more training steps.
This is called domain randomisation~\cite{tobin2017dr,tan2018sim2real}.
Teacher--student distillation~\cite{lee2020terrain,miki2022perceptive} allow to reduce the size of the final network for online inference.

\subsection{Approach comparison}
\label{sec:bg-axes}

The main difference stem between the RL approach and the rest.
Previous approach yielded a fixed mathematical law that could be analysed with the classical control toolkit.
The RL approach is more like a black box where most of the classical tool for stability analysis do not apply.


\subsection{Reinforcement learning for locomotion}
\label{sec:bg-rl-loco}

\subsubsection{Control architecture}

Several actuator-level control architectures have been used by the literature.

\textbf{Direct torque control} provides the policy with the highest level of control authority, but also imposes the most demanding communication and control-loop requirements. It is the mode the quasi-direct-drive actuator lineage was designed for~\cite{seok2015design,wensing2017proprioceptive}, and the one MIT Cheetah commands through~\cite{dicarlo2018mpc}. In simulation, the policy can directly command joint torques at a high frequency. On the physical robot, however, closing such a fast loop requires the policy, communication bus, actuator controller, and state estimation to operate with sufficiently low latency. In the present hardware architecture, this approach was not practical through the CAN communication bus at the required frequency.

\textbf{Direct position control} substantially reduces these requirements by delegating the low-level actuator dynamics to the motor controller. Position targets over a joint-level PD loop are what most learned locomotion uses~\cite{hwangbo2019anymal,tan2018sim2real,rudin2021}. Two variants were investigated: conventional smart-servo position control and the MIT control mode with fixed compliance parameters~\cite{katz2019minicheetah}. Both allow the RL policy to operate at a lower frequency while the actuator controller handles the fast inner loop.

\textbf{Impedance control} provides an intermediate level of control authority. Instead of directly specifying the actuator torque or position response, the policy can influence the desired mechanical behaviour through position, stiffness, damping, and/or torque-related parameters. The $[q_d, \dot q_d, K_p, K_d, \tau_{\text{ff}}]$ command frame established by the Cheetah lineage is the standard realisation~\cite{katz2019minicheetah,wensing2017proprioceptive}, and \cref{eq:impedance} is the form this thesis uses. This allows the low-level controller to preserve fast and stabilising actuator dynamics while giving the policy greater freedom than pure position control.

\subsubsection{Gait prior}

The second major design dimension is the amount and location of prior structure imposed on the gait.

A strong prior can be introduced through a \textbf{reference trajectory}. In this case, the policy is primarily trained to track or adapt an authored motion. DeepMimic tracks motion-capture clips this way~\cite{peng2018deepmimic}, and the Cassie line tracked authored reference motions~\cite{xie2019}. This strongly restricts the reachable set of motions and can substantially simplify learning. However, the quality of the resulting behaviour is bounded by the choice of reference: there is no guarantee that the selected trajectory is optimal for the particular robot morphology and dynamics.

A weaker prior can instead be introduced through a \textbf{trajectory generator} whose parameters are modulated by the policy~\cite{pmtg2018}. This preserves the assumption that locomotion should have a structured, periodic form while allowing the learned controller to search for a gait adapted to the plant. Trajectory generators can themselves provide different levels of freedom.

At the stronger end of this generator-based formulation, the policy only modifies a small number of \textbf{high-level gait parameters}, such as phase, frequency, stride length, or amplitude~\cite{pmtg2018,cpgrl2022}. Central pattern generators are the classical form of that prior~\cite{ijspeert2008}, and dynamical movement primitives the classical form for non-periodic motions~\cite{ijspeert2013dmp}. The shape of the individual joint trajectories remains largely prescribed. This provides a relatively small search space but limits the set of gaits that can be represented.

A weaker prior allows the policy to modify the \textbf{complete trajectory}. Rather than predicting every trajectory sample independently, the possible trajectories can be represented using a low-dimensional basis, such as Fourier series~\cite{shafii2010tfs,yang2010tfs,wang2023fourier,fld2024} or Bézier curves~\cite{westervelt2007book,castillo2020hzdrl}. The policy then modifies the coefficients of this representation. This retains useful structural assumptions, most importantly periodicity and smoothness, while allowing the shape of the gait itself to become a learned variable.

At the extreme, \textbf{no explicit gait prior} is imposed and the policy directly learns the complete joint trajectory~\cite{heess2017emergence,radosavovic2024humanoid}. Periodic locomotion can still be encouraged through reward shaping, but periodicity then becomes an emergent property rather than a constraint. Such an approach provides a large reachable set but considerably increases the exploration and learning problem, and there is no guarantee that the learned solution will converge to a periodic gait.

This distinction is important because the periodicity prior is not unique to the present framework. It can be introduced at different levels: through the \textbf{reward}~\cite{siekmann2021}, through a \textbf{reference motion}~\cite{peng2018deepmimic}, through an authored \textbf{trajectory generator}~\cite{pmtg2018,cpgrl2022}, or directly through the \textbf{action space} by parameterising the trajectory itself~\cite{fld2024,ates2025freqgait}. These choices determine how much of the gait is fixed and how much is available to the policy at run time.

\subsubsection{Choice for DASH-01}

For DASH-01, a fully unconstrained joint-level policy was considered undesirable because it would require reinforcement learning to simultaneously discover both the appropriate periodic structure of the gait and the stabilising behaviour of the robot. Conversely, imposing a fixed reference trajectory would strongly constrain the solution and assume that a sufficiently good gait was already known.

DASH-01 also does not correspond closely to a particular biological morphology. Consequently, transferring a reference trajectory derived from human or animal locomotion would introduce a potentially unjustified assumption about the optimal gait. A trajectory generator therefore provides a useful intermediate solution: the existence of a periodic gait is used as a prior, while its detailed shape remains a decision variable for the policy.

The framework consequently investigates a \textbf{Fourier-based gait parameterisation}, in which the joint trajectories over a gait cycle are represented by Fourier coefficients. The policy can modify these coefficients to search for a gait adapted to the specific morphology and dynamics of DASH-01. This reduces the dimensionality and imposes smooth periodic structure compared with a fully free trajectory representation, while retaining substantially more freedom than tracking a fixed reference trajectory.

The use of a Fourier representation should not itself be interpreted as the contribution of this work. Fourier-based periodic trajectory representations and the use of basis coefficients for locomotion predate this framework. The contribution of the framework instead lies in how the gait parameterisation is integrated with reinforcement learning and fast feedback control, and in the experimental comparison of alternative parameterisations under equivalent conditions.

In particular, the framework distinguishes between the \textbf{slow gait-generation channel} and the \textbf{fast feedback channel}. The gait parameters determine the underlying periodic motion, while the feedback component is responsible for reacting to disturbances and short-timescale deviations from that nominal gait. This separation allows the learned policy to modify the gait without requiring it to reconstruct the complete high-frequency actuator behaviour at every timestep.

The choice of this architecture was therefore motivated by a trade-off between \textbf{reachable gait space and learning complexity}. A fixed reference trajectory provides a strong prior but may exclude the optimal gait. A completely free policy provides a very weak prior but greatly increases the search space. A parameterised trajectory generator occupies an intermediate position: it constrains the search to smooth periodic motions while allowing the gait itself to be discovered by reinforcement learning.

\subsubsection{Architectures investigated}

Because the original objective of the thesis was to develop an RL framework capable of producing high-speed bipedal locomotion, these architectural choices were treated as experimental design variables rather than assumptions fixed from the outset.

Multiple actuator and gait-control architectures were therefore implemented and evaluated during framework development. These included direct position control, MIT-mode control with fixed compliance parameters, impedance-based control, gait-generator approaches with different levels of policy authority, and architectures incorporating a hand-tuned reflex controller alongside the learned component.

The final architecture was selected based on the resulting trade-off between actuator controllability, communication and bandwidth requirements, exploration complexity, and the freedom available to discover a gait adapted to DASH-01. The Fourier gait representation should consequently be understood as the selected design point within this broader architectural investigation, rather than as a novel gait representation in itself.


A \textbf{shape prior} fixes the trajectory and exposes its timing, so tuning does not enlarge the
reachable set. A \textbf{structural prior} constrains the class and leaves the member free, so the
reachable set grows with the parameterisation. DASH-01 frees the intra-limb axis because it has no
template to copy. Hardware calibration established that the cam is a full crank, sweeping about
$245^{\circ}$ against a CAD-guessed $\pm86^{\circ}$ (\cref{sec:hw-experiments}), and the 4-bar's
transmission ratio varies strongly with posture (\cref{sec:hw-force-map}). No animal offers a gait
for that morphology. The mirror on the inter-limb axis is a shape prior, taken because the task is
a straight-line dash and because it halves the emitted specification to 14 coefficients. Its cost
was paid later: a mirror-symmetric gait travels only in a straight line, so steering had to be
added as two explicit asymmetry channels.

Three trade-offs follow from the lineage, and the thesis lives on all three.

\textbf{Sim-to-real gap.} Ranked as measured on this robot rather than as usually listed, the
sources are actuator dynamics and drive bandwidth, mass and inertia error, the contact and friction
model, latency, sensor noise, and unmodelled compliance. The standard toolkit is domain
randomisation with curricula, observation history in place of privileged state, explicit sensor and
delay models, and actuator networks~\cite{hwangbo2019anymal}. The useful framing is randomisation
as robustness training over a distribution of plants, rather than as noise injection. This thesis
measures first, and randomises around the measurement (\cref{sec:sim2real-lineage}). The
counter-lesson is that randomising around an optimistic plant does not recover the plant error. No
policy has run on the hardware, so that is an inference from the plant-correction table rather than
a transfer result.

\textbf{Learning efficiency and dimensionality reduction.} The raw problem is a torque sequence in
$\mathbb{R}^{n_{\text{joint}} \times T}$, and the structured problem is 10 to 50 gait-morphology
parameters. The prior asserted is that successful locomotion is approximately periodic,
band-limited and bounded. The reduction removes exploration rather than imposing answers. It costs
a reachable set bounded by the harmonic count, no stability guarantee from smoothness, and the need
for a fast feedback channel outside the generator. Modern results use $10^8$ to $10^9$ transitions,
with thousands of parallel environments~\cite{rudin2021}. A curriculum optimal at that budget is
not necessarily optimal at this project's, so sample efficiency is the actual research variable
here.

\textbf{Curriculum design.} Published curricula range over command, terrain, disturbance and
randomisation, usually on an already fully dynamic robot, and the recent shift is toward
performance-adaptive schedules. This thesis instead applies difficulty to the
\emph{controllability of the plant} (\cref{sec:rl-curriculum}), by releasing base degrees of
freedom one at a time. It is weaker than the state of the art on every other axis: binary rather
than continuous release, hand-scheduled rather than performance-gated advance, and a single
difficulty axis rather than a vector. The known danger is that a strongly constrained plant lets
the policy learn that a suppressed mode does not exist, so release then produces immediate failure.

\subsection{Platforms, and where DASH-01 sits}
\label{sec:bg-platforms}

The bipeds producing the results above fall into three groups, separated more by actuation and
budget than by control philosophy. \textbf{Series-elastic research bipeds} such as
ATRIAS~\cite{hubicki2016atrias} and its descendants Cassie and Digit carry leaf springs in series
with the drives, and a high-voltage bus. They hold the speed record~\cite{crowley2024cassie}, and
they are the product of a decade of institutional effort. \textbf{Quasi-direct-drive robots in the
MIT Cheetah lineage}~\cite{seok2015design,wensing2017proprioceptive,katz2019minicheetah} are
predominantly quadrupeds, and architecturally the most influential group. Their pattern is
quasi-direct-drive actuation commanded in torque mode, adopted by essentially every recent legged
platform. \textbf{Hobby-class humanoids} are built from position-controlled smart servos, and are
confined to the static-walking regime for a structural reason rather than a budget one. A servo
closing its own position loop cannot deliver a commanded ground force, so the dynamic families stay
unavailable whatever controller sits above them.

DASH-01 sits deliberately between the second and third groups: quasi-direct-drive actuators and a
running-oriented morphology, on a student budget and a single-board computer. Three bets define it.
A parallel 4-bar knee moves the knee actuator's mass proximally. A passive ankle spring replaces a
seventh and eighth motor. Six motors do the work of ten. \Cref{tab:speed-ceiling} confirms the
first bet, at $-0.93$\,(m/s)/kg for distal mass against $-0.11$ for the torso.
\Cref{sec:hw-force-map} shows the second bet was wrong by a factor of eight. The third stands.

What separates this platform from its neighbours is not any single choice. It is that the drive
electronics were bought as a bundle with the actuator. That put the robot in the third group's
actuation interface, while its mechanics belong to the second. \Cref{sec:plant-id-drive} measures
what it costs.
